% ============================================================================
%  Chapitre 4 — Analyse des besoins, conception et mise en œuvre
%  Le travail ne consiste pas à concevoir une chaîne ex nihilo mais à
%  diagnostiquer, faire évoluer et outiller une chaîne CI/CD existante.
% ============================================================================
\chapter{Analyse des besoins, conception et mise en œuvre des évolutions de la chaîne CI/CD}
\label{chap:besoins}

\section{Analyse de l'existant et recueil des besoins}
\label{sec:bes-analyse}

\subsection{Situation à mon arrivée}

À mon arrivée dans l'équipe COM, trois workflows GitHub Actions étaient en
place : le premier compilait la couche applicative, exécutait les tests
unitaires et l'analyse statique Coverity ; le deuxième construisait la
distribution Linux embarquée avec Yocto ; le troisième produisait l'image de
l'environnement de compilation. Cette organisation offrait un premier niveau
d'automatisation, mais restait perfectible sur le suivi des indicateurs de
qualité, la centralisation des résultats, l'observabilité des chaînes et
l'automatisation des contrôles de sécurité. Ma mission consistait donc à
reprendre cette chaîne, à en identifier les limites et à la faire évoluer,
non à en concevoir une nouvelle.

\subsection{Méthode de recueil}

La prise en main a combiné trois sources. La documentation interne
(Confluence) et l'étude des dépôts GitHub Enterprise (workflows, scripts,
configuration, historique des modifications) ont permis de reconstituer les
étapes de construction et les dépendances entre dépôts. Des échanges réguliers
avec trois développeurs de l'équipe COM et le Product Owner ont confronté cette
documentation au fonctionnement réel et fait émerger les attentes des
utilisateurs. Enfin, mon tuteur, coordinateur DevSecOps, m'a aidé à distinguer
ce qui était imposé par l'environnement de l'entreprise de ce qui pouvait être
adapté au module de communication. Ce recueil s'est poursuivi de manière
itérative tout au long du stage.

\subsection{Besoins identifiés}

L'analyse n'a révélé aucun dysfonctionnement bloquant : l'enjeu était
d'accompagner la croissance des besoins du projet. Les besoins identifiés se
regroupent en quatre axes :

\begin{itemize}[nosep]
  \item \textbf{Reproductibilité} : enrichir, optimiser et versionner l'image de
        compilation, et y intégrer le \gls{sdk} de la cible pour disposer d'un
        environnement commun à la compilation native et croisée ;
  \item \textbf{Automatisation} : prendre en charge plusieurs configurations
        (compilation native AMD64 pour les tests et la couverture, compilation
        croisée AArch64 pour le produit) et fiabiliser les scripts qui lancent
        les analyses, attendent leur fin, récupèrent les résultats et produisent
        les rapports ;
  \item \textbf{Qualité et sécurité} : mesurer la couverture des tests (lignes,
        fonctions, branches) et analyser la composition binaire de l'artefact
        livré avec BDBA pour identifier les composants tiers et leurs
        vulnérabilités ;
  \item \textbf{Observabilité} : dépasser la consultation d'une exécution
        isolée en collectant, historisant et restituant les indicateurs dans des
        tableaux de bord Grafana.
\end{itemize}

% ===========================================================================
\section{Formalisation des exigences}
\label{sec:bes-exigences}

\subsection{Parties prenantes}

Les développeurs COM sont les utilisateurs directs de la chaîne ; le Product
Owner en exploite les résultats pour le suivi du projet ; le coordinateur
DevSecOps veille à la cohérence avec les pratiques de l'entreprise. Les services
internes (exécuteurs, serveurs d'analyse, référentiel d'artefacts) n'expriment
pas de besoin fonctionnel, mais imposent des règles d'utilisation
(tableau~\ref{tab:parties-prenantes}).

\begin{table}[H]
\centering
\caption{Parties prenantes et attentes principales}
\label{tab:parties-prenantes}
\small
\begin{tabularx}{\textwidth}{@{}p{3.4cm} X@{}}
\toprule
\textbf{Partie prenante} & \textbf{Attentes principales} \\
\midrule
Développeurs COM & Résultats rapides et lisibles sur la compilation, les tests,
  la couverture et les analyses de sécurité \\
Product Owner & Visibilité sur l'avancement, l'évolution de la qualité et les
  principaux risques \\
Coordinateur DevSecOps & Chaîne automatisée, maintenable, sécurisée et
  compatible avec les services de l'entreprise \\
Services internes & Respect des politiques de sécurité, des restrictions d'accès
  et des conditions d'utilisation \\
\bottomrule
\end{tabularx}
\end{table}

\subsection{Exigences fonctionnelles}

Les exigences fonctionnelles (tableau~\ref{tab:bes-ef}) sont regroupées par
chaîne et pour la solution d'observabilité.

\begin{table}[H]
\centering
\caption{Exigences fonctionnelles}
\label{tab:bes-ef}
\small
\begin{tabularx}{\textwidth}{@{}l X@{}}
\toprule
\textbf{Réf.} & \textbf{La solution doit\ldots} \\
\midrule
\multicolumn{2}{@{}l}{\emph{Chaîne applicative}} \\
EF1 & compiler l'application en natif AMD64 et en croisé AArch64 \\
EF2 & prendre en charge les modes \texttt{DEBUG} et \texttt{RELEASE} \\
EF3 & exécuter l'analyse statique et récupérer les défauts détectés \\
EF4 & soumettre l'artefact cible à l'analyse de composition binaire \\
EF5 & exécuter les tests unitaires et produire un rapport \\
EF6 & exécuter les tests d'intégration compatibles avec l'intégration continue \\
EF7 & mesurer la couverture du code par les tests unitaires \\
EF8 & publier une synthèse lisible des compilations, tests, couverture et analyses \\
EF9 & produire un artefact rattaché à la révision, à la configuration et à l'environnement \\
\midrule
\multicolumn{2}{@{}l}{\emph{Chaîne de la distribution}} \\
EF10 & construire automatiquement la distribution Linux embarquée \\
EF11 & produire les variantes d'image, notamment de développement et de production \\
EF12 & générer le \gls{sdk} utilisé pour la compilation croisée \\
EF13 & publier les images et le \gls{sdk} dans le référentiel d'artefacts \\
\midrule
\multicolumn{2}{@{}l}{\emph{Chaîne de l'image de compilation}} \\
EF14 & construire une image contenant compilateurs, bibliothèques, outils d'analyse et \gls{sdk} \\
EF15 & vérifier la présence et les versions des principaux outils \\
EF16 & valider l'image candidate par une construction applicative complète \\
EF17 & versionner et publier l'image validée dans un registre \\
\midrule
\multicolumn{2}{@{}l}{\emph{Observabilité}} \\
EF18 & collecter automatiquement les indicateurs produits par les chaînes \\
EF19 & transformer les indicateurs dans un format commun \\
EF20 & historiser les données pour les comparer dans le temps \\
EF21 & restituer les indicateurs dans des tableaux de bord Grafana \\
EF22 & visualiser l'évolution des constructions, tests, couverture, défauts et vulnérabilités \\
\bottomrule
\end{tabularx}
\end{table}

\subsection{Exigences non fonctionnelles}

Les exigences non fonctionnelles (tableau~\ref{tab:bes-enf}) précisent les
propriétés attendues de la solution.

\begin{table}[H]
\centering
\caption{Exigences non fonctionnelles}
\label{tab:bes-enf}
\small
\begin{tabularx}{\textwidth}{@{}l X@{}}
\toprule
\textbf{Réf.} & \textbf{Exigence} \\
\midrule
\multicolumn{2}{@{}l}{\emph{Reproductibilité}} \\
ENF1 & Versions des outils de compilation et d'analyse épinglées dans l'image \\
ENF2 & Même image utilisée sur les postes et dans la chaîne d'intégration \\
ENF3 & Artefacts rattachés à la révision, à la configuration et à la version de l'environnement \\
\midrule
\multicolumn{2}{@{}l}{\emph{Performance}} \\
ENF4 & Durée des chaînes compatible avec le rythme de développement \\
ENF5 & Configurations indépendantes exécutées en parallèle \\
ENF6 & Caches partagés pour la construction de la distribution \\
ENF7 & Optimisations sans perte de fiabilité ni de traçabilité \\
\midrule
\multicolumn{2}{@{}l}{\emph{Maintenabilité et réutilisabilité}} \\
ENF8 & Logique placée dans des scripts indépendants de l'orchestrateur \\
ENF9 & Workflows paramétrables et réutilisables \\
ENF10 & Évolutions documentées pour l'équipe \\
ENF11 & Composants génériques facilitant une mutualisation COM/MET \\
\midrule
\multicolumn{2}{@{}l}{\emph{Sécurité}} \\
ENF12 & Aucun secret en clair dans les dépôts ni dans les couches de l'image \\
ENF13 & Jetons injectés à l'exécution, de portée minimale \\
ENF14 & Aucun secret exposé dans les journaux d'exécution \\
\midrule
\multicolumn{2}{@{}l}{\emph{Observabilité}} \\
ENF15 & Données des différents outils transformées dans un format homogène \\
ENF16 & Collecte automatique, sans intervention manuelle \\
ENF17 & Indicateurs définis et documentés permettant de suivre la maturité du module \\
ENF18 & Tableaux de bord lisibles et exploitables par l'équipe \\
\midrule
\multicolumn{2}{@{}l}{\emph{Infrastructure d'exécution}} \\
ENF19 & Exécution sur des machines éphémères sans état local \\
ENF20 & Dépendances aux chemins et caches locaux limitées ou externalisées \\
ENF21 & Compatibilité avec les politiques réseau et de sécurité de l'entreprise \\
\bottomrule
\end{tabularx}
\end{table}

Deux sujets sont volontairement hors de ces exigences : l'apport de
l'intelligence artificielle au diagnostic des échecs, traité comme un axe
exploratoire, et la migration vers une chaîne commune aux équipes COM et MET,
que la solution doit seulement faciliter (ENF11).

\subsection{Contraintes et hypothèses}

La plateforme de code source, l'orchestrateur, les ressources d'exécution et
plusieurs outils d'analyse sont imposés par l'entreprise. Le travail consiste
donc à adapter la chaîne à ces services, non à choisir une plateforme.

\paragraph{Infrastructure d'exécution.}
\label{sssec:bes-contraintes-infra}
Les traitements s'exécutent sur des machines éphémères, détruites à la fin de
chaque job. L'isolation est meilleure, mais aucun fichier, dépendance ou
résultat intermédiaire n'est conservé : les constructions longues exigent un
cache partagé, indépendant de la durée de vie des exécuteurs. Les fonctions de
GitHub Actions qui supposent une machine persistante ou un accès privilégié ne
sont pas disponibles. Les chemins de travail varient, ce qui impose de les
normaliser pour les outils qui corrèlent les défauts par chemin de fichier.
Enfin, des règles encadrent les droits d'exécution, l'accès réseau et
l'usage des secrets.

\paragraph{Distribution Linux embarquée.} Sa construction mobilise un grand
nombre de composants et dure longtemps sans réutilisation des résultats
intermédiaires ; sa maîtrise dépend des caches de téléchargement et d'état
partagé et d'un parallélisme adapté aux exécuteurs. Ses livrables (images et
\gls{sdk}) évoluent à un rythme différent du code applicatif, ce qui justifie
une chaîne séparée.

\paragraph{Outils d'entreprise.} Les outils sont imposés ; la marge de décision
porte sur l'organisation et le paramétrage des workflows, la répartition des
jobs, les scripts exécutables en local, la configuration des outils, la
sécurisation des accès, la gestion des caches, la restitution des résultats et
la collecte des indicateurs.

\paragraph{Hypothèses.} Les dépôts, les services internes et les secrets sont
disponibles depuis les exécuteurs autorisés ; les tests sur cible matérielle
restent hors de la chaîne automatique ; les images et le \gls{sdk} produits par
la chaîne de distribution servent de référence pour la compilation croisée ;
les indicateurs collectés sont représentatifs de la qualité du firmware.

% ===========================================================================
\section{Conception des évolutions de la chaîne CI/CD}
\label{sec:bes-conception}

\subsection{Principes directeurs}
\label{ssec:bes-principes}

Six principes ont guidé la conception des évolutions. Ils découlent directement
des limites identifiées et des contraintes de l'infrastructure d'exécution.

\begin{description}[style=nextline,leftmargin=0.8cm]
  \item[P1 --- La logique réside dans des scripts, pas dans le YAML] Les
    fichiers de workflow se limitent à l'orchestration (déclencheurs, matrices,
    dépendances entre jobs, injection des secrets). Tout traitement non trivial
    est placé dans le dépôt \texttt{bluebird/ci}, sous \texttt{scripts/} :
    \texttt{cov\_build.sh} (phases capture, analyse et commit de Coverity),
    \texttt{extract\_cov.sh} et \texttt{pyCovXtract.py} (extraction des
    défauts), \texttt{gcov\_summary\_md.sh} et \texttt{robot\_summary\_md.py}
    (synthèses Markdown), \texttt{BDBA\_GetReport.py} et
    \texttt{BDBA\_ExtractMetrics.py} (analyse de composition). Les scripts
    \texttt{build\_pm77xx-app.sh} et \texttt{build\_pm77xx-yocto.sh} servent de
    références pour rejouer une construction en local. Ce principe répond à
    ENF8 : un script se teste sur un poste, un fichier de workflow ne se teste
    qu'en poussant un commit.
  \item[P2 --- Un seul environnement de compilation] La même image de
    conteneur, \texttt{bbbuild\_cov\_ut}, sert à la compilation native amd64 et
    à la compilation croisée aarch64. Elle embarque la chaîne native (GCC
    13.3.0, CMake 3.28.3), Coverity 2024.12.0 et le \gls{sdk} de la cible
    installé dans \texttt{/usr/local/oecore-x86\_64}. Le basculement vers la
    cible se limite à sourcer le script \texttt{environment-setup-aarch64-evo-linux}
    du \gls{sdk} avant l'appel à CMake (ENF1, ENF2).
  \item[P3 --- Ne publier une image qu'après l'avoir éprouvée] Une image
    candidate est publiée sous l'étiquette \texttt{:test}, validée en lui faisant
    construire l'application complète par la chaîne \texttt{pm77xx-app}, puis
    seulement promue sous son étiquette définitive (EF16).
  \item[P4 --- Des paramètres explicites, avec des valeurs par défaut sûres]
    Chaque entrée du workflow est reprise dans une variable d'environnement
    \texttt{INPUT\_*} munie d'une valeur par défaut, par exemple
    \texttt{INPUT\_APP\_REPO\_BRANCH: "\$\{\{ inputs.app\_repo\_branch || 'develop' \}\}"}.
    Les exécutions planifiées, pour lesquelles le contexte \texttt{inputs} est
    vide, se comportent ainsi comme une exécution manuelle aux valeurs par
    défaut (ENF9).
  \item[P5 --- Les secrets n'existent que le temps de l'exécution] La clé de
    service Git, les jetons Coverity, BDBA, Artifactory et VictoriaMetrics sont
    fournis par le mécanisme \texttt{secrets} de la plateforme. Les fichiers
    d'authentification sont écrits avec les droits \texttt{0600}, le fichier de
    connexion Coverity (\texttt{.cov\_cnx.yml}) est supprimé dès la fin de
    l'extraction, et la construction de l'image utilise un montage SSH éphémère
    (ENF12 à ENF14).
  \item[P6 --- Restituer le résultat là où le développeur regarde déjà] Chaque
    job écrit une synthèse dans \texttt{\$GITHUB\_STEP\_SUMMARY} (tests,
    couverture, défauts, composants vulnérables). Le développeur lit son résultat
    dans la page du run, sans ouvrir Coverity, BDBA ou Grafana (EF8).
\end{description}

\subsection{Contraintes imposées et décisions d'intégration}
\label{ssec:bes-decisions}

Le tableau~\ref{tab:bes-decisions} distingue, pour chaque brique, ce qui est
imposé par l'environnement de l'entreprise et la décision d'intégration prise.
La colonne de droite reprend les points qui relèvent de mon travail.

\begin{table}[H]
\centering
\caption{Éléments imposés et décisions d'intégration}
\label{tab:bes-decisions}
\small
\begin{tabularx}{\textwidth}{@{}p{2.3cm} X X@{}}
\toprule
\textbf{Brique} & \textbf{Imposé} & \textbf{Décision d'intégration} \\
\midrule
Orchestration & GitHub Enterprise Actions ; exécuteurs éphémères
  \texttt{orchestrator-runner} (service RaaS) ; conteneur lancé avec
  \texttt{-u root} &
  Trois workflows distincts ; matrice DEBUG/RELEASE $\times$ amd64/aarch64 ;
  \texttt{workflow\_call} pour réutiliser la chaîne applicative ; jobs légers
  (rapports, BDBA, publication) déportés sur \texttt{ubuntu-22.04} \\
\addlinespace
Analyse statique & Coverity 2024.12.0 ; serveur Coverity interne, avec un
  projet et un flux dédiés &
  Trois phases séparées ; build instrumenté limité à DEBUG/aarch64 ;
  \texttt{--strip-path} pour stabiliser les chemins ; vues P1 et P2
  extraites par \texttt{pyCovXtract.py} \\
\addlinespace
Composition & Black Duck Binary Analysis (BDBA) &
  Analyse de l'archive RELEASE/aarch64 uniquement ; scan par
  \texttt{Foxy/action-bdba-scan} puis récupération du rapport et des métriques
  par deux scripts Python via l'API \\
\addlinespace
Artefacts & Instance Artifactory de l'entreprise ; un dépôt générique et un
  registre d'images &
  Arborescence \texttt{sdp-artifacts/}, \texttt{sdk/}, \texttt{bsp-images/} ;
  nommage \texttt{<nom>-<version>-<horodatage>} ; publication limitée à la
  branche \texttt{main} pour l'application \\
\addlinespace
Distribution & Yocto, branche \texttt{scarthgap} du dépôt
  \texttt{evo-image-config} &
  Pilotage par kas 5.2 plutôt que par BitBake direct ; une cellule de matrice
  par variante d'image ; génération du \gls{sdk} sur la variante de
  développement \\
\addlinespace
Métriques & Instance VictoriaMetrics interne ;
  authentification Basic &
  Action composite \texttt{upload-vm-metrics} ; deux formats d'import selon la
  source ; étiquettes \texttt{bu}, \texttt{project}, \texttt{repo},
  \texttt{branch}, \texttt{workflow} \\
\bottomrule
\end{tabularx}
\end{table}

Deux décisions méritent d'être explicitées.

\paragraph{Où placer l'instrumentation Coverity.} Le build Coverity et le build
de test sont deux constructions distinctes. Le build capturé par
\texttt{cov-build} cible l'architecture livrée (aarch64) : les unités de
compilation analysées correspondent au code qui part en production. Les tests
et la couverture s'exécutent sur un build amd64 non instrumenté par Coverity, ce
qui évite de superposer l'instrumentation \texttt{cov-build} et
l'instrumentation \texttt{gcov}. Dans le workflow, les deux étapes de
compilation sont mutuellement exclusives par une condition \texttt{if} portée
par la matrice.

\paragraph{Intégrer le \gls{sdk} dans l'image de compilation.} Le
\gls{sdk} est produit par la chaîne Yocto, publié dans
Artifactory, puis installé dans l'image de compilation par le
\texttt{Dockerfile}. Les deux chaînes sont ainsi couplées par un artefact
versionné et non par un état partagé.

\subsection{Architecture cible}
\label{ssec:bes-architecture}

\begin{figure}[H]
\centering
\includegraphics[width=\textwidth]{workflow-diagram}
\caption{Architecture générale de la chaîne d'intégration et de livraison}
\label{fig:bes-architecture}
\end{figure}

La figure~\ref{fig:bes-architecture} montre le flux. Un déclencheur (planifié
en jours ouvrés, manuel ou appel depuis un autre workflow) démarre un exécuteur
éphémère. Le job récupère l'image de compilation depuis Artifactory, puis, dans
ce conteneur, extrait trois dépôts : \texttt{bluebird/ci} (scripts),
\texttt{pm77xx-app/sdp} (code applicatif, sous-modules inclus) et
\texttt{pm77xx-app/automation-tests} (suites Robot Framework). Les cellules de
la matrice compilent, analysent et testent en parallèle ; les jobs aval
téléchargent l'archive produite, génèrent les rapports, soumettent l'archive de
production à BDBA, la publient dans Artifactory et envoient les métriques à
VictoriaMetrics, lues par Grafana.

Le tableau~\ref{tab:bes-herite-evolue} précise la part héritée et la part
apportée.

\begin{table}[H]
\centering
\caption{Situation à mon arrivée et évolutions apportées}
\label{tab:bes-herite-evolue}
\small
\begin{tabularx}{\textwidth}{@{}p{3.3cm} X X@{}}
\toprule
\textbf{Domaine} & \textbf{Existant} & \textbf{Évolution apportée} \\
\midrule
Compilation & Compilation de l'application &
  Matrice DEBUG/RELEASE $\times$ native amd64 / croisée aarch64 \\
Environnement & Image de compilation &
  Intégration du \gls{sdk} ; optimisation du \texttt{Dockerfile} ; mise à jour des paquets \\
Qualité & Analyse statique Coverity ; tests unitaires &
  Couverture (\texttt{gcovr}) et synthèse ; tests d'intégration Robot Framework \\
Sécurité & --- & Scan BDBA, rapport PDF et métriques \\
Livraison & Construction de la distribution &
  Publication des artefacts applicatifs, des images et du \gls{sdk} \\
Observabilité & Résultats par exécution &
  Action de collecte vers VictoriaMetrics ; mise à jour du tableau de bord Grafana \\
Infrastructure & Exécuteurs éphémères &
  Étude des solutions NFS et snapshot ; correction de scripts bash et Python \\
\bottomrule
\end{tabularx}
\end{table}

\subsection{Décomposition en chaînes}
\label{ssec:bes-chaines}

Le système comporte trois chaînes aux rythmes et aux durées différents
(tableau~\ref{tab:bes-chaines}). Le dépôt contient également un workflow
\texttt{evo-bench-image}, destiné aux images du banc de test ; il est hors du
périmètre de ce mémoire.

\begin{table}[H]
\centering
\caption{Chaînes d'intégration}
\label{tab:bes-chaines}
\small
\begin{tabularx}{\textwidth}{@{}l X p{4cm}@{}}
\toprule
\textbf{Chaîne} & \textbf{Objet} & \textbf{Déclenchement} \\
\midrule
\texttt{pm77xx-app} & Compilation, analyses, tests et publication de la couche applicative &
  Cron \texttt{0 0 * * 1-5}, manuel, appelé \\
\texttt{pm77xx-yocto} & Cinq variantes d'image et \gls{sdk} &
  Cron \texttt{0 22 * * 1-5}, manuel \\
\texttt{publish-image} & Construction, validation et publication de l'image de compilation &
  Manuel \\
\bottomrule
\end{tabularx}
\end{table}

La chaîne applicative est planifiée à minuit UTC et la chaîne Yocto à 22 h UTC :
elles ne se disputent donc pas les exécuteurs. Les durées mesurées sont
présentées au chapitre~\ref{chap:evaluation}.

La séparation s'impose pour trois raisons. Premièrement, les rythmes diffèrent :
le code applicatif change tous les jours alors que la distribution et l'image
de compilation évoluent rarement. Deuxièmement, les durées sont incompatibles :
intégrer une construction Yocto à plus d'une heure dans le chemin du développeur
dégraderait le retour d'information (ENF4). Troisièmement, les chaînes
s'articulent par des artefacts versionnés : le \gls{sdk} sort de
\texttt{pm77xx-yocto}, entre dans \texttt{publish-image}, et
\texttt{publish-image} appelle \texttt{pm77xx-app} comme sous-workflow pour
valider l'image candidate.

\begin{figure}[H]
\centering
\includegraphics[width=\textwidth]{vue-global}
\caption{Historique des exécutions de \texttt{pm77xx-app} depuis la création du workflow (1\,186 exécutions)}
\label{fig:bes-historique}
\end{figure}

% ===========================================================================
\section{Mise en œuvre technique}
\label{sec:impl}

Chaque sous-section suit la même trame : l'objectif visé, l'état de départ,
l'évolution apportée, puis la difficulté rencontrée. Les références EF et ENF
renvoient aux exigences de la section précédente.

\subsection{Les chaînes CI}
\label{ssec:impl-chaines}

\subsubsection{Chaîne applicative \texttt{pm77xx-app}}
\label{sec:impl-pipeline-app}

\emph{Exigences : EF1 à EF9, ENF4, ENF5, ENF8, ENF9.}

\paragraph{Paramètres.} La chaîne accepte les entrées du
tableau~\ref{tab:impl-parametres-app} par \texttt{workflow\_dispatch} et
\texttt{workflow\_call}. Le mode appelé existe pour la validation de l'image de
compilation : \texttt{publish-image} appelle \texttt{pm77xx-app} en lui
passant l'étiquette candidate \texttt{bbbuild\_cov\_ut:test}, avec
\texttt{run\_cov\_analysis=false} pour éviter de polluer le serveur Coverity avec
une image non validée.

\begin{table}[H]
\centering
\caption{Paramètres d'entrée de \texttt{pm77xx-app}}
\label{tab:impl-parametres-app}
\small
\begin{tabularx}{\textwidth}{@{}l X l@{}}
\toprule
\textbf{Paramètre} & \textbf{Rôle} & \textbf{Défaut} \\
\midrule
\texttt{app\_repo\_branch} & Branche de \texttt{pm77xx-app/sdp} & \texttt{develop} \\
\texttt{test\_repo\_branch} & Branche de \texttt{automation-tests} & \texttt{main} \\
\texttt{builder\_img\_tag} & Image de compilation & \texttt{bbbuild\_cov\_ut:2.3} \\
\texttt{ci\_repo\_branch} & Branche de \texttt{bluebird/ci} & \texttt{main} \\
\texttt{run\_cov\_analysis} & Active la capture Coverity & \texttt{true} \\
\texttt{run\_robot\_tests} & Active les tests d'intégration & \texttt{false} \\
\texttt{workspace\_path} & Répertoire de travail dans le conteneur & \texttt{/\_\_w/ci/ci} \\
\bottomrule
\end{tabularx}
\end{table}

\paragraph{Matrice de construction.} La construction est déclinée selon
\texttt{build\_type} $\in$ \{DEBUG, RELEASE\} et \texttt{compiler\_type}
$\in$ \{amd64, aarch64\}, soit quatre cellules indépendantes
(\texttt{fail-fast: false}), exécutées en parallèle. Le
tableau~\ref{tab:impl-matrice-activites} répartit les activités sur ces
cellules ; cette répartition est un arbitrage entre coût et couverture.

\begin{table}[H]
\centering
\caption{Répartition des activités sur la matrice}
\label{tab:impl-matrice-activites}
\small
\begin{tabular}{@{}lcccc@{}}
\toprule
\textbf{Activité} & \makecell{DEBUG\\amd64} & \makecell{DEBUG\\aarch64} & \makecell{RELEASE\\amd64} & \makecell{RELEASE\\aarch64} \\
\midrule
Compilation            & \checkmark & \checkmark & \checkmark & \checkmark \\
Analyse statique       & --- & \checkmark & --- & --- \\
Tests unitaires        & \checkmark & --- & \checkmark & --- \\
Couverture             & \checkmark & --- & --- & --- \\
Tests d'intégration    & \checkmark$^{*}$ & --- & --- & --- \\
Analyse de composition & --- & --- & --- & \checkmark \\
Publication            & --- & --- & --- & \checkmark$^{\dagger}$ \\
\bottomrule
\end{tabular}

\smallskip
{\footnotesize $^{*}$ si \texttt{run\_robot\_tests=true}. $^{\dagger}$ sur la branche \texttt{main} uniquement.}
\end{table}

Les tests s'exécutent sur amd64 parce que le binaire s'exécute nativement sur
l'exécuteur ; les cellules aarch64 ne peuvent que compiler. La couverture
exige l'instrumentation \texttt{gcov}, donc un build DEBUG. L'analyse de
composition porte sur RELEASE/aarch64, c'est-à-dire sur l'artefact réellement
livré. Le listing~\ref{lst:impl-build} montre le point de bascule entre
compilation native et croisée.

\begin{lstlisting}[language=yaml,caption={Compilation native ou croisée selon la cellule de matrice},label={lst:impl-build}]
- name: Build ${{ matrix.compiler_type }} App in ${{ matrix.build_type }} mode
  id: build_app
  if: ${{ env.INPUT_RUN_COV_ANALYSIS == 'false'
       || matrix.build_type != 'DEBUG'
       || matrix.compiler_type != 'aarch64' }}
  run: |
    if [ "${{ matrix.compiler_type }}" == 'aarch64' ]; then
      . /usr/local/oecore-x86_64/environment-setup-aarch64-evo-linux
    fi
    cmake -DCMAKE_BUILD_TYPE=${{ matrix.build_type }} \
      -S "${{ env.SDP_ROOT }}" \
      -B "${{ env.SDP_ROOT }}/${{ matrix.build_type }}-build.${{ matrix.compiler_type }}"
    cmake --build "..." --parallel
\end{lstlisting}

La condition \texttt{if} rend cette étape exclusive de l'étape
\texttt{cov\_build\_app}, qui porte la condition inverse. Le répertoire de
build est suffixé par la configuration (\texttt{DEBUG-build.aarch64}), ce qui
évite toute collision entre cellules.

\begin{figure}[H]
\centering
\includegraphics[width=0.95\textwidth]{workflow-pm77xx-app}
\caption{Exécution de \texttt{pm77xx-app} (\#1316) : matrice de construction et jobs aval}
\label{fig:impl-activite-app}
\end{figure}

La figure~\ref{fig:impl-activite-app} se lit de gauche à droite. Le job
\emph{Show Input Parameters} (0 s) écrit les paramètres dans la synthèse. Les
quatre cellules de la matrice démarrent ensuite : la cellule DEBUG/amd64 dure
19 min 28 s, contre 4 min 55 s à 5 min 17 s pour les trois autres ; elle
cumule la compilation instrumentée, les tests, la couverture et, si activés,
les tests d'intégration, et constitue donc le chemin critique des 23 min 06 s
de l'exécution. Les jobs aval (couverture 26 s, BDBA 56 s, publication 25 s,
rapports de test 28 s et 24 s) sont déclenchés dès la fin de la matrice
(\texttt{needs: build-scan-test}).

\paragraph{Difficulté : clonage multi-dépôts.} Le code est réparti entre un
dépôt principal et de nombreux sous-modules privés. Le checkout utilise
\texttt{submodules: recursive} avec la clé de service \texttt{GIT\_SRV\_SSH\_KEY}
et \texttt{persist-credentials: false}, de sorte que la clé ne reste pas dans
la configuration Git du répertoire de travail.

\subsubsection{Tests unitaires et couverture}
\label{ssec:impl-ut}

\emph{Exigences : EF5, EF7, EF8.}

\paragraph{Tests unitaires.} Les tests sont lancés par \texttt{ctest}, avec
export JUnit pour la publication, et excluent la suite \texttt{mcu\_unit\_tests},
destinée au cœur temps réel et inexécutable sur l'hôte :

\begin{lstlisting}[language=bash,caption={Exécution des tests unitaires},label={lst:impl-ctest}]
SDP_ROOT="$SDP_ROOT" ctest \
  --test-dir "$SDP_ROOT/${BUILD_TYPE}-build.${ARCH}" \
  --exclude-regex "mcu_unit_tests" \
  --output-on-failure \
  --output-junit "$SDP_ROOT/${BUILD_TYPE}-build.${ARCH}/test_results.xml"
\end{lstlisting}

L'étape est bornée par \texttt{timeout-minutes: 1} : un test qui ne rend jamais
la main ne doit pas immobiliser l'exécuteur jusqu'à la limite globale. Le job \texttt{test-reports} publie ensuite les
résultats avec l'action \texttt{EnricoMi/publish-unit-test-result-action}
(13 tests réussis sur 13 dans l'exécution \#1316).

\paragraph{Couverture.} La couverture est mesurée avec \texttt{gcovr} sur la
seule cellule DEBUG/amd64 :

\begin{lstlisting}[language=bash,caption={Mesure de couverture},label={lst:impl-gcovr}]
gcovr --root "$SDP_ROOT" --object-directory . \
  --exclude '.*/test/.*' \
  --gcov-ignore-parse-errors=negative_hits.warn \
  --html-details coverage/gcov.html \
  --lcov coverage/gcov.info \
  --json coverage/gcov.json \
  --json-summary coverage/gcov-summary.json \
  --print-summary | tee coverage/gcov-summary.txt
\end{lstlisting}

Trois choix sont à relever. L'option
\texttt{--gcov-ignore-parse-errors=negative\_hits.warn} dégrade en avertissement
les compteurs d'exécution négatifs produits par \texttt{gcov}, qui, sans cette
option, interrompent la génération du rapport. Les quatre formats de sortie répondent à quatre usages : HTML pour la
lecture, LCOV pour les outils tiers, JSON pour l'exploitation, texte pour la
synthèse. Enfin, le pourcentage de lignes est extrait par un filtre
\texttt{awk} sur la ligne \texttt{lines:} du résumé pour alimenter
VictoriaMetrics.

\texttt{gcov\_summary\_md.sh} convertit le résumé en tableau Markdown avec
badge et statut par métrique, selon deux seuils (50 \% et 75 \%). Sur le run
\#1316, la couverture est de 32,3 \% des lignes (19\,308 / 59\,772), 20,9 \% des
fonctions (1\,511 / 7\,216) et 29,6 \% des branches (5\,559 / 18\,775) ; les trois
indicateurs sont sous le seuil bas. La chaîne rend donc visible un niveau de
couverture faible plutôt que de le masquer.

\begin{figure}[H]
\centering
\includegraphics[width=0.75\textwidth]{summary-pm77xx-app}
\caption{Synthèse d'exécution : défauts Coverity, tests, couverture et BDBA}
\label{fig:impl-synthese}
\end{figure}

\paragraph{Difficulté : chemins d'artefacts.} \texttt{upload-artifact} conserve
l'arborescence absolue du conteneur. Le job \texttt{ut-coverage-reports}, qui
s'exécute dans une autre image, doit donc retrouver les fichiers sous
\texttt{./build-amd64/\_\_w/ci/ci/pm77xx/pm77xx-app/sdp/DEBUG-build.amd64}.
Ce chemin est centralisé dans une variable \texttt{BUILD\_PATH} pour limiter sa
duplication.

\subsubsection{Tests d'intégration Robot Framework}
\label{ssec:impl-robot}

\emph{Exigence : EF6.}

Les suites Robot Framework du dépôt \texttt{automation-tests} exercent les
services applicatifs. Leur exécution est un ajout récent. Elle est désactivée par
défaut (\texttt{run\_robot\_tests=false}) et ne concerne que la cellule
DEBUG/amd64, car le lanceur de services attend les binaires de l'arborescence
x86\_64/DEBUG. La séquence, entièrement conditionnée à \texttt{steps.launch\_sdp.outcome}, est la suivante :

\begin{enumerate}[nosep]
  \item \texttt{cmake --install} copie configuration, scripts de migration de base
        de données et certificats sous \texttt{/opt/app-ext/sdp} ;
  \item \texttt{sdp.launch.in.container.sh} démarre les services en arrière-plan
        (\texttt{timeout-minutes: 5}) ;
  \item \texttt{robot} produit \texttt{output.xml}, \texttt{log.html},
        \texttt{report.html} et \texttt{xunit.xml} ;
  \item \texttt{sdp.launch.in.container.sh --kill} arrête les services dans une
        étape \texttt{if: always()}, même si les tests ont échoué ;
  \item \texttt{robot\_summary\_md.py} écrit la synthèse (badge de taux de
        réussite, compteurs, 50 échecs listés au plus) dans le résumé du run.
\end{enumerate}

La difficulté principale est d'exécuter un ensemble de services système
(\texttt{dbus}, services dépendant de \texttt{libsystemd}) dans un conteneur sans
gestionnaire de services complet. Le lanceur démarre donc chaque service
lui-même et la garde de 5 minutes protège contre un service qui resterait au
premier plan.

\subsubsection{Analyse statique Coverity}
\label{ssec:impl-coverity}

\emph{Exigence : EF3.}

Coverity préexistait ; mon travail a porté sur la fiabilisation des scripts et
sur la restitution. L'analyse suit trois phases pilotées par
\texttt{cov\_build.sh \{capture|analyze|commit\}}, chacune dans une étape
distincte afin d'isoler les échecs et de mesurer chaque durée. La configuration
\texttt{scripts/coverity\_ci.yaml} active tous les vérificateurs, y compris la
famille de sécurité C :

\begin{lstlisting}[language=yaml,caption={Configuration Coverity},label={lst:impl-cov-config}]
capture:
  build:
    cov-build-args: [--parallel-translate=4]
analyze:
  c-cpp-virtual: true
  c-cpp-fnptr: true
  checkers: { all: true, rule: true, c-family-security: true }
  cov-analyze-args:
    - --strip-path
    - /__w/ci/ci/pm77xx/pm77xx-app/
commit:
  connect:
    stream: "<flux-projet>"
    url: "https://<serveur-coverity>"
\end{lstlisting}

\texttt{c-cpp-virtual} et \texttt{c-cpp-fnptr} résolvent les appels virtuels
et les pointeurs de fonction, nombreux dans le code applicatif.
\texttt{--strip-path} supprime le préfixe du répertoire de travail afin que
l'identité d'un défaut (fichier, fonction, type) reste stable d'un run à l'autre
malgré des exécuteurs éphémères (ENF20). Sans lui, un même défaut serait
compté comme nouveau à chaque modification du chemin.

\paragraph{Restitution ciblée.} \texttt{extract\_cov.sh} interroge, pour chacune
des vues P1 et P2, l'API REST
\texttt{/api/v2/streams/viewContents/\{viewid\}} via \texttt{pyCovXtract.py},
génère les rapports \texttt{coverity\_report\_<run>\_P1.html} et
\texttt{\_P2.html}, et expose \texttt{cov\_p1\_issue\_nbr},
\texttt{cov\_p2\_issue\_nbr} et \texttt{cov\_total\_issue\_nbr} dans
\texttt{GITHUB\_OUTPUT}. Sur le run \#1316 : 3 défauts P1 et 46 défauts P2.

\paragraph{Politique de blocage.} Le script retourne 1 lorsque des défauts sont
trouvés. L'étape l'invoque avec \texttt{|| true}, car un code de sortie non nul
colorait l'étape en rouge alors que le résultat est déjà exposé dans la
synthèse. Seuls un échec de capture, d'analyse ou de connexion fait donc échouer
le job. Ce choix est assumé : avec une dette initiale de 46 défauts P2, un
blocage immédiat paralyserait l'équipe. La trajectoire cible est un blocage sur
les défauts nouvellement introduits uniquement, ce que l'option \texttt{--diff}
de \texttt{pyCovXtract.py} permettrait de calculer.

\subsubsection{Analyse de composition BDBA}
\label{ssec:impl-bdba}

\emph{Exigence : EF4.}

Le job \texttt{scan-bdba} télécharge l'archive RELEASE/aarch64 puis enchaîne
trois étapes :

\begin{enumerate}[nosep]
  \item soumission par \texttt{Foxy/action-bdba-scan@v1.1.1}
        (\texttt{delete-binary: true}, produit \texttt{pm77xx-app} version 1.0.0) ;
  \item \texttt{BDBA\_GetReport.py} interroge \texttt{/api/product/\{id\}/} toutes
        les 30 s jusqu'au statut \texttt{R} (délai maximal de 600 s) puis télécharge
        le rapport PDF ;
  \item \texttt{BDBA\_ExtractMetrics.py} produit un fichier JSON Lines au format
        d'import VictoriaMetrics, et expose les compteurs
        \texttt{bdba\_components\_total}, \texttt{bdba\_components\_vuln},
        \texttt{bdba\_vulns\_total}, \texttt{bdba\_vulns\_critical},
        \texttt{bdba\_vulns\_high}, \texttt{bdba\_vulns\_medium},
        \texttt{bdba\_vulns\_low}.
\end{enumerate}

Le tableau~\ref{tab:impl-bdba-metriques} définit les indicateurs. Sur
l'exécution capturée (figure~\ref{fig:impl-bdba}), BDBA identifie 4 composants
dont 3 vulnérables (\texttt{cjson}, \texttt{openldap}, \texttt{hostapd}) pour
122 vulnérabilités : 5 critiques, 19 élevées, 86 moyennes et 12 faibles. Les
compteurs BDBA de la synthèse (figure~\ref{fig:impl-synthese}) proviennent d'une
autre exécution et diffèrent donc légèrement ; ce mémoire retient les valeurs
du rapport.

\begin{table}[H]
\centering
\caption{Indicateurs extraits de l'analyse de composition}
\label{tab:impl-bdba-metriques}
\small
\begin{tabularx}{\textwidth}{@{}l X@{}}
\toprule
\textbf{Indicateur} & \textbf{Définition} \\
\midrule
\texttt{bdba\_component} & Composants tiers détectés, étiquetés \texttt{vulnerable} ou \texttt{clean} \\
\texttt{bdba\_vulns} & Vulnérabilités connues, étiquetées par sévérité \\
\bottomrule
\end{tabularx}
\end{table}

\begin{figure}[H]
\centering
\includegraphics[width=0.7\textwidth]{rapport-bdba}
\caption{Rapport BDBA de l'archive \texttt{pm77xx-app} (RELEASE/aarch64)}
\label{fig:impl-bdba}
\end{figure}

Cette analyse alimente l'inventaire des composants que les obligations de
nomenclature logicielle exigent (section~\ref{sec:art-supply-chain}).

\subsubsection{Publication des artefacts applicatifs}
\label{ssec:impl-publication-app}

\emph{Exigence : EF9.}

Chaque cellule empaquette son résultat avant de le publier :
\texttt{DESTDIR=/tmp/staging cmake --install} installe dans une arborescence
temporaire, puis \texttt{tar -czf} archive \texttt{lib} et \texttt{bin} depuis
\texttt{/opt/app-ext/sdp}. L'archive est conservée comme artefact de run
(\texttt{<BUILD>-build-<ARCH>}). Pour la branche \texttt{main}, le job
\texttt{upload-artifact} publie l'archive RELEASE/aarch64 dans Artifactory par
\texttt{jf rt upload} dans le répertoire \texttt{sdp-artifacts/} du dépôt
générique, sous un nom de la forme
\texttt{pm77xx-app-artifact-<version>-<horodatage>.tar.gz}.
La taille de l'archive est aussi mesurée (\texttt{du -cb}) et envoyée comme
métrique \texttt{sdp\_artifact\_size} (78,0 Mo sur la capture du tableau de bord).

\emph{Limite :} la version \texttt{ARTIFACT\_VERSION} est fixée à \texttt{0.0.1} ;
seul l'horodatage distingue deux artefacts, et le nom ne contient ni
l'empreinte du commit ni la version de l'image de compilation. L'exigence EF9
n'est donc que partiellement couverte.

\subsubsection{Chaîne de la distribution \texttt{pm77xx-yocto}}
\label{sec:impl-pipeline-yocto}

\emph{Exigences : EF10 à EF13, ENF6.}

La configuration est décrite avec kas 5.2 en trois niveaux, composés dans la
ligne de commande par des « : » :

\begin{lstlisting}[language=bash,caption={Appels kas : image et SDK},label={lst:impl-kas}]
kas build kas/board/<carte>.yml:kas/project/image-evo-dev.yml
kas build kas/board/<carte>.yml:kas/project/image-evo-dev.yml:\
         kas/project/option/populate-sdk.yml
\end{lstlisting}

Le premier niveau décrit la carte (\Carte), le second
le projet d'image, le troisième des options, dont \texttt{populate-sdk.yml}.
Pour la variante de production, le fichier de carte est suffixé
(\texttt{<carte>-prod.yml}) grâce à l'entrée \texttt{include} de la matrice.
Les cinq variantes (tableau~\ref{tab:impl-variantes}) sont construites en
parallèle.

\begin{table}[H]
\centering
\caption{Variantes d'image produites}
\label{tab:impl-variantes}
\small
\begin{tabularx}{\textwidth}{@{}l X@{}}
\toprule
\textbf{Variante} & \textbf{Usage} \\
\midrule
\texttt{image-evo-dev} & Développement, outils de débogage ; sert aussi à générer le \gls{sdk} \\
\texttt{image-evo-hwtest} & Validation matérielle en production \\
\texttt{image-evo-netboot} & Démarrage réseau (initramfs) \\
\texttt{image-evo-rescue} & Image de secours (initramfs) \\
\texttt{image-evo-prod} & Production \\
\bottomrule
\end{tabularx}
\end{table}

\begin{figure}[H]
\centering
\includegraphics[width=0.6\textwidth]{workflow-pm77xx-yocto}
\caption{Exécution de \texttt{pm77xx-yocto} (\#543, 1 h 05 min 27 s)}
\label{fig:impl-yocto}
\end{figure}

\paragraph{Difficulté : BitBake refuse de s'exécuter en root.} Le conteneur
est imposé en root par l'infrastructure (section~\ref{ssec:impl-runners}), or
BitBake refuse de s'exécuter sous cet utilisateur. Le workflow écrit donc la clé
SSH sous \texttt{/home/user/.ssh}, transfère la propriété du répertoire de
travail à \texttt{user} (\texttt{chown -R}), déclare le dépôt en
\texttt{safe.directory} puis lance \texttt{kas} par \texttt{su user -c}.

\paragraph{Publication.} Après la construction, le \gls{sdk} est localisé par
\texttt{find build/*tmp*/deploy/sdk -name '*-toolchain-*.sh'} puis publié dans
\texttt{sdk/} ; les variantes \texttt{dev} et \texttt{prod} sont archivées
(\texttt{find -type f | tar --transform 's|.*/||'} met tous les fichiers à la racine de
l'archive) et publiées dans \texttt{bsp-images/}. Pour la variante dev, une
archive « flashable » regroupe \texttt{deploy/images} et le fichier machine
\texttt{<carte>.conf} dans \texttt{bsp-images/standalone-flashing-images/}.
Un développeur applicatif obtient ainsi un compilateur croisé et une image de
référence sans reconstruire la distribution.

\begin{figure}[H]
\centering
\includegraphics[width=0.75\textwidth]{jfrog-artifactory}
\caption{Organisation du dépôt générique Artifactory}
\label{fig:impl-artifactory}
\end{figure}

\subsubsection{Chaîne de l'image de compilation \texttt{publish-image}}
\label{sec:impl-conteneur}

\emph{Exigences : EF14 à EF17, ENF1, ENF2, ENF12.}

\paragraph{Construction du \texttt{Dockerfile}.} L'image est construite en trois
étages à partir d'une image Ubuntu 24.04 (\texttt{se-ubuntu-noble:3.0.1}) :
\texttt{compiler} (paquets et compilateurs), \texttt{cov-sdk\_builder}
(Coverity et \gls{sdk}) et \texttt{app\_builder} (image finale). Seul le contenu
utile est copié dans l'étage final (\texttt{COPY --from=cov-sdk\_builder}), de sorte que l'archive
Coverity et l'installateur du \gls{sdk} n'alourdissent pas l'image.

\begin{table}[H]
\centering
\caption{Principaux composants de l'image et versions épinglées}
\label{tab:impl-image-versions}
\small
\begin{tabularx}{\textwidth}{@{}l X l@{}}
\toprule
\textbf{Composant} & \textbf{Rôle} & \textbf{Version} \\
\midrule
GCC / G++ & Compilation C et C++ & 13.3.0-6ubuntu2\textasciitilde 24.04.1 \\
CMake & Système de construction & 3.28.3-1build7 \\
Meson / Ninja & Brique de cybersécurité (csb-app) & 1.3.2 / 1.11.1 \\
Coverity Analysis & Analyse statique & 2024.12.0 \\
GoogleTest / GMock & Tests unitaires & 1.14.0-1 \\
CMocka & Doublures de dépendances & 1.1.7-3 \\
gcovr & Couverture & 7.0-1 \\
Doxygen / Graphviz & Documentation & 1.9.8 / 2.42.2 \\
Python & Scripts d'intégration continue & 3.12.3 \\
\bottomrule
\end{tabularx}
\end{table}

Les optimisations apportées sont les suivantes :

\begin{itemize}[nosep]
  \item versions épinglées par paquet (\texttt{cmake=3.28.3-1build7}), pour que deux
        constructions séparées dans le temps utilisent la même chaîne d'outils (ENF1) ;
  \item \texttt{--no-install-recommends}, qui évite l'installation de paquets non demandés ;
  \item caches BuildKit \texttt{--mount=type=cache,target=/var/cache/apt} et
        \texttt{/var/lib/apt} (\texttt{sharing=locked}) : les paquets déjà téléchargés
        sont réutilisés d'une construction à l'autre sans rester dans une couche ;
  \item construction multi-étages, décrite plus haut.
\end{itemize}

\paragraph{Accès aux dépôts privés.} La brique de cybersécurité \texttt{csb-app}
(version \texttt{Release\_csbApp\_2.6.0}) est clonée et compilée avec Meson dans
l'image. Le clonage utilise \texttt{RUN --mount=type=ssh}, alimenté par
\texttt{docker build --ssh default} : la clé reste dans l'agent SSH de l'hôte et
n'est écrite dans aucune couche de l'image.

\begin{lstlisting}[language=docker,caption={Installation du SDK et clonage sécurisé},label={lst:impl-dockerfile}]
COPY ./resources/*sdk-installer*.sh ${SDK_PATH}/sdk-installer.sh
RUN chmod +x sdk-installer.sh \
    && ./sdk-installer.sh -d /usr/local/oecore-x86_64 -y

RUN --mount=type=ssh \
    ssh-keyscan <forge-interne> >> /root/.ssh/known_hosts \
    && git clone --branch Release_csbApp_2.6.0 \
       git@<forge-interne>:<organisation>/csb-app.git csb-app
\end{lstlisting}

L'image définit l'utilisateur \texttt{user}, mais son exécution en
intégration se fait en root (\texttt{options: -u root}) pour des raisons de
droits d'écriture imposées par RaaS ; ce point est reporté en perspectives.

\begin{figure}[H]
\centering
\includegraphics[width=\textwidth]{workflow-image-docker}
\caption{Exécution de \texttt{publish-image} : construction, contrôles, validation, publication}
\label{fig:impl-cycle-image}
\end{figure}

\paragraph{Cycle de validation.} La figure~\ref{fig:impl-cycle-image} détaille
l'enchaînement :
\begin{enumerate}[nosep]
  \item \texttt{build\_builder\_image} (11 min 08 s) construit avec
        \texttt{DOCKER\_BUILDKIT=1 docker build --ssh default --build-arg COV\_VERSION=2024.12.0}
        et pousse l'image \texttt{:test} ;
  \item \texttt{lint\_dockerfile} exécute \texttt{hadolint}, et \texttt{Check Image}
        (3 min 20 s) lance \texttt{display\_version.sh} dans l'image pour vérifier la
        présence et les versions des outils (EF15) ;
  \item \emph{Image Validation} appelle \texttt{pm77xx-app} avec
        \texttt{builder\_img\_tag=bbbuild\_cov\_ut:test} : quatre cellules de
        compilation, rapports de test, couverture et BDBA ;
  \item \emph{Image Publication} (28 min 09 s) revérifie l'étiquette, republie
        l'image sous son étiquette définitive et en exporte une archive.
\end{enumerate}

Une image n'atteint donc l'étiquette définitive qu'après avoir prouvé qu'elle
construit le produit. Sur la capture, le job \texttt{lint\_dockerfile} est en
échec alors que l'exécution se termine en succès : les erreurs de \texttt{hadolint}
sont tolérées. C'est un choix provisoire, le temps de résorber les
avertissements existants.

% ===========================================================================
\subsection{L'infrastructure d'exécution}
\label{ssec:impl-runners}

\emph{Exigences : ENF4, ENF6, ENF19 à ENF21.}

\subsubsection{Modèle d'exécuteurs éphémères}

Les jobs lourds (compilation, Yocto, construction d'image) s'exécutent sur
\texttt{orchestrator-runner}, un service d'exécuteurs qui alloue une machine
neuve par job et la détruit ensuite. Les jobs légers (rapports de test, couverture, BDBA,
publication) utilisent \texttt{ubuntu-22.04} avec l'image
\texttt{se-ubuntu-noble:3.0.1} : ils n'ont pas besoin de la chaîne d'outils et
ne doivent pas occuper un exécuteur coûteux. Les contraintes générales de ce
modèle sont décrites en section~\ref{sssec:bes-contraintes-infra} ; leur
traduction dans la chaîne est la suivante :

\begin{itemize}[nosep]
  \item \textbf{Droits.} Le conteneur est lancé avec \texttt{options: -u root},
        ce qui oblige la chaîne Yocto à rétrograder les privilèges
        (section~\ref{sec:impl-pipeline-yocto}).
  \item \textbf{Chemins.} Le répertoire de travail \texttt{/\_\_w/ci/ci} est
        paramétré (\texttt{workspace\_path}) et neutralisé côté Coverity par
        \texttt{--strip-path}.
  \item \textbf{Temps.} Des garde-fous \texttt{timeout-minutes} limitent les étapes
        susceptibles de bloquer.
\end{itemize}

\subsubsection{Limite : le cache de la distribution}

Les répertoires \texttt{DL\_DIR} et \texttt{SSTATE\_DIR} sont créés dans
l'espace de travail du job par un fichier \texttt{site.conf} généré à
l'exécution :

\begin{lstlisting}[language=bash,caption={Configuration Yocto générée par le job},label={lst:impl-siteconf}]
DL_DIR     = '$GITHUB_WORKSPACE/projets/yocto/downloads'
SSTATE_DIR = '$GITHUB_WORKSPACE/projets/yocto/sstate-cache'
BB_NUMBER_THREADS = '32'
PARALLEL_MAKE = '-j 32'
\end{lstlisting}

L'exécuteur étant détruit à la fin du job, ces deux caches repartent de zéro à
chaque exécution : chacune des cinq variantes retélécharge les sources et
recompile l'ensemble des paquets, ce qui explique en grande partie la durée de
1 h 05 de la chaîne. La parallélisation à 32 threads ne réduit que la durée de
chaque compilation, pas la redondance entre cellules. L'exigence ENF6
(cache partagé) n'est donc \emph{pas} couverte par la chaîne actuelle.

\subsubsection{Étude des solutions NFS et snapshot}
\label{sssec:impl-etude-cache}

J'ai étudié deux solutions pour rendre ces caches persistants malgré le caractère
éphémère des exécuteurs (tableau~\ref{tab:impl-cache}).

\begin{table}[H]
\centering
\caption{Solutions étudiées pour le cache de la distribution}
\label{tab:impl-cache}
\small
\begin{tabularx}{\textwidth}{@{}p{2.6cm} X X@{}}
\toprule
\textbf{Solution} & \textbf{Principe} & \textbf{Points d'attention} \\
\midrule
Partage NFS &
  Montage d'un volume réseau comme \texttt{DL\_DIR} et \texttt{SSTATE\_DIR} &
  Écritures concurrentes des cinq cellules de la matrice ; latence sur de très
  nombreux petits fichiers ; disponibilité du service depuis les exécuteurs \\
Snapshot &
  Volume restauré au démarrage de l'exécuteur à partir d'un instantané du cache &
  Instantané à rafraîchir régulièrement ; cache en lecture seule, les nouveaux
  états n'étant pas réinjectés ; compatibilité avec le service d'exécuteurs \\
\bottomrule
\end{tabularx}
\end{table}

\paragraph{Conclusion de l'étude.} La solution retenue est le partage NFS. Le
snapshot ne fournit qu'un cache en lecture seule, figé à la date de
l'instantané : les états produits par une exécution ne profitent pas aux
suivantes tant que l'instantané n'est pas régénéré. Un volume NFS, au contraire,
est alimenté à chaque construction et reste à jour sans opération
supplémentaire. Le principal point à traiter lors de la mise en œuvre est la
gestion des écritures concurrentes des cinq cellules de la matrice. Cette
solution n'est pas encore mise en œuvre ; elle constitue la première perspective
de ce travail (section~\ref{sec:concl-perspectives}).

\subsubsection{Gestion des secrets et sécurité}

Les secrets utilisés sont \texttt{GIT\_SRV\_SSH\_KEY}, \texttt{COV\_AUTH\_JSON},
\texttt{COV\_SRV\_CNXON}, \texttt{AM\_BDBA\_API\_TOKEN}, \texttt{JFROG\_SRV\_PAT} et
\texttt{VMET\_PASSWORD}. Les fichiers qui les contiennent sont créés avec
\texttt{chmod 0600} et supprimés après usage (\texttt{.cov\_cnx.yml}), le
checkout n'enregistre pas les identifiants (\texttt{persist-credentials: false})
et l'image de compilation ne contient aucune clé grâce au montage SSH. Deux
limites subsistent : le jeton BDBA est passé en argument de ligne de commande
aux scripts Python, donc visible dans la liste des processus de l'exécuteur
(masqué dans les journaux, mais pas dans la machine) ; et la chaîne Yocto écrit
la clé SSH sur le disque de l'exécuteur, dans un répertoire détruit avec lui.

% ===========================================================================
\subsection{L'observabilité}
\label{ssec:impl-observabilite}

\emph{Exigences : EF18 à EF22, ENF15 à ENF18.}

La synthèse d'un run ne montre pas de tendance. Pour suivre la qualité dans le
temps, chaque chaîne envoie ses indicateurs à VictoriaMetrics, qu'un tableau de
bord Grafana interroge.

\paragraph{Action de collecte.} L'action composite
\texttt{.github/actions/upload-vm-metrics} encapsule l'envoi pour éviter de
dupliquer les appels \texttt{curl} dans chaque workflow. Elle accepte deux
modes, selon la source :

\begin{lstlisting}[language=bash,caption={Deux modes d'envoi des métriques},label={lst:impl-vm}]
# Mode fichier : JSON Lines d'import (BDBA, plusieurs séries)
curl -T "$METRICS_JSONL_FILE" -u "$VMET_USER:$VMET_PASSWORD" \
     -X POST "https://$VMET_URL/api/v1/import"

# Mode valeur unique : format texte Prometheus
curl -d "${METRIC_NAME}{${METRIC_LABELS}} ${METRIC_VALUE}" \
     -u "$VMET_USER:$VMET_PASSWORD" \
     -X POST "https://$VMET_URL/api/v1/import/prometheus"
\end{lstlisting}

Le mode texte suffit pour une valeur scalaire (taille de l'artefact,
couverture) ; le mode JSON Lines est nécessaire quand un script produit
plusieurs séries étiquetées (\texttt{bdba\_component} par statut,
\texttt{bdba\_vulns} par sévérité). Les étiquettes communes
(\texttt{bu="em"}, \texttt{project="bluebird"}, \texttt{repo},
\texttt{branch}, \texttt{workflow}) permettent de filtrer par chaîne et par
branche dans Grafana. Le nom d'utilisateur est une variable de dépôt et le mot de
passe un secret.

\begin{table}[H]
\centering
\caption{Indicateurs collectés vers VictoriaMetrics}
\label{tab:impl-indicateurs}
\small
\begin{tabularx}{\textwidth}{@{}l X l@{}}
\toprule
\textbf{Métrique} & \textbf{Contenu} & \textbf{Origine} \\
\midrule
\texttt{sdp\_artifact\_size} & Taille de l'archive \texttt{lib} + \texttt{bin} & \texttt{du -cb}, chaque cellule \\
\texttt{code\_coverage} & Couverture de lignes (\%) & \texttt{gcovr}, DEBUG/amd64 \\
\texttt{bdba\_component} & Composants, par statut & \texttt{BDBA\_ExtractMetrics.py} \\
\texttt{bdba\_vulns} & Vulnérabilités, par sévérité & \texttt{BDBA\_ExtractMetrics.py} \\
\bottomrule
\end{tabularx}
\end{table}

\paragraph{Tableau de bord.} J'ai mis à jour le tableau de bord Grafana
« Bluebird w1.2 / Daily -- COM team » (figure~\ref{fig:impl-grafana}). Il est
organisé en quatre sections : sprint en cours, \emph{Quality \& Security}
(défauts Coverity P1 et P2, couverture, taille de l'artefact, répartition des
composants et des vulnérabilités), état des workflows (statut et durée du
dernier run, un onglet par chaîne) et pull requests ouvertes. À la date de la
capture, il affiche 3 défauts P1, 46 défauts P2, 32,3 \% de couverture, un
artefact de 78,0 Mo et un dernier run \texttt{pm77xx-app} réussi en 23,1 minutes.

\begin{figure}[H]
\centering
\includegraphics[width=0.6\textwidth]{dashboard-grafana}
\caption{Tableau de bord Grafana de l'équipe COM}
\label{fig:impl-grafana}
\end{figure}

% ===========================================================================
\subsection{Bilan}
\label{ssec:impl-bilan}

\subsubsection{Couverture des exigences}

Le tableau~\ref{tab:impl-bilan-ef} indique, pour chaque groupe d'exigences
fonctionnelles, la section qui les traite et leur état.

\begin{table}[H]
\centering
\caption{Couverture des exigences fonctionnelles}
\label{tab:impl-bilan-ef}
\small
\begin{tabularx}{\textwidth}{@{}l X l l@{}}
\toprule
\textbf{Réf.} & \textbf{Constat} & \textbf{Section} & \textbf{État} \\
\midrule
EF1, EF2 & Matrice native/croisée $\times$ DEBUG/RELEASE & \ref{sec:impl-pipeline-app} & Couverte \\
EF3 & Coverity en trois phases, rapports P1/P2 & \ref{ssec:impl-coverity} & Couverte \\
EF4 & Scan BDBA de l'archive RELEASE/aarch64 & \ref{ssec:impl-bdba} & Couverte \\
EF5, EF7, EF8 & Tests unitaires, couverture, synthèse & \ref{ssec:impl-ut} & Couverte \\
EF6 & Robot Framework, désactivé par défaut, DEBUG/amd64 & \ref{ssec:impl-robot} & Partielle \\
EF9 & Version fixe, pas d'empreinte de commit dans le nom & \ref{ssec:impl-publication-app} & Partielle \\
EF10 à EF13 & Cinq variantes, \gls{sdk}, publication & \ref{sec:impl-pipeline-yocto} & Couverte \\
EF14 à EF17 & Image versionnée, validée, publiée & \ref{sec:impl-conteneur} & Couverte \\
EF18, EF20, EF21 & Collecte, historisation, Grafana & \ref{ssec:impl-observabilite} & Couverte \\
EF19, EF22 & Deux formats d'import ; évolution des constructions limitée au dernier run & \ref{ssec:impl-observabilite} & Partielle \\
\bottomrule
\end{tabularx}
\end{table}

Pour les exigences non fonctionnelles, ENF1, ENF2, ENF4, ENF5, ENF8, ENF9,
ENF12 à ENF14 et ENF19 à ENF21 sont couvertes par les choix décrits plus haut.
ENF3 (traçabilité des artefacts) est partielle, pour la raison donnée pour
EF9. ENF6 (cache partagé de la distribution) n'est pas couverte : l'étude
NFS/snapshot a conduit à retenir la solution NFS. ENF11 (mutualisation entre les équipes COM et
MET) n'a pas été traitée.

\subsubsection{Difficultés transverses}

Le débogage des scripts bash et Python a occupé une part notable du travail.
Les problèmes récurrents étaient :

\begin{itemize}[nosep]
  \item un code de sortie significatif (\texttt{extract\_cov.sh}) interprété comme un
        échec par la plateforme ;
  \item des chemins qui varient selon l'image dans laquelle s'exécute le job
        (artefacts, répertoire de travail) ;
  \item des étapes qui ne terminent jamais, d'où les \texttt{timeout-minutes} ;
  \item des citations de secrets et de chaînes dans les scripts embarqués dans le YAML,
        déplacés dans \texttt{scripts/} pour être testables (principe P1).
\end{itemize}

Les limites qui découlent de ce bilan sont discutées en
section~\ref{sec:eval-discussion}, et les perspectives correspondantes en
section~\ref{sec:concl-perspectives}.

